\chapter{ТЕХНИЧЕСКАЯ РЕАЛИЗАЦИЯ ЭКСПЕРИМЕНТАЛЬНОГО ФРЕЙМВОРКА}

% -------------------------------------------------------
\section{Архитектура воспроизводимого экспериментального фреймворка}
% -------------------------------------------------------

\subsection{Общая структура кодовой базы}

Экспериментальный фреймворк реализован в виде единого Python-проекта с общей инфраструктурой и тремя независимыми экспериментальными модулями. Управление зависимостями осуществляется через менеджер пакетов \texttt{uv} с фиксированным lockfile, что обеспечивает побайтовую воспроизводимость окружения. Минимальная требуемая версия интерпретатора~--- Python~3.12.

Структура проекта:
\begin{itemize}
  \item \texttt{src/common/} --- общая инфраструктура: асинхронный клиент OpenRouter (\texttt{openrouter.py}), утилита загрузки данных из Google Sheets (\texttt{google\_sheets.py}), учёт потреблённых токенов;
  \item \texttt{common\_config.py} --- централизованная конфигурация: базовый URL API, список тестируемых моделей, загрузка ключа из переменной окружения \texttt{OPENROUTER\_API\_KEY};
  \item \texttt{linda-problem/} --- эксперимент 1: \texttt{run\_experiments.py}, \texttt{compute\_metrics.py}, \texttt{generate\_dataset.py}, \texttt{config.py};
  \item \texttt{wason-selection/} --- эксперимент 2: \texttt{run\_experiments.py}, \texttt{compute\_metrics.py}, \texttt{download\_inputs.py}, \texttt{build\_graphs.py}, \texttt{upload\_results.py};
  \item \texttt{wason-2-4-6/} --- эксперимент 3: \texttt{run\_experiments.py}, \texttt{metrics.py}, \texttt{compute\_metrics.py}, \texttt{build\_graphs.py};
  \item \texttt{tests/} --- unit-тесты для всех трёх экспериментальных модулей.
\end{itemize}

Основные зависимости: \texttt{openai} (клиент OpenRouter), \texttt{pandas}, \texttt{numpy}, \texttt{scipy}, \texttt{scikit-learn} (статистический анализ), \texttt{matplotlib}, \texttt{seaborn} (визуализация), \texttt{pydantic} (валидация конфигураций), \texttt{google-auth}, \texttt{gspread} (интеграция с Google Sheets).

Все три эксперимента запускаются через единый интерфейс командной строки:
\begin{itemize}
  \item эксперимент 1: \texttt{uv run python linda-problem/run\_experiments.py {-}{-}api-key \$OPENROUTER\_API\_KEY}
  \item эксперимент 2: \texttt{uv run python wason-selection/run\_experiments.py {-}{-}models <model\_id> {-}{-}api-key \$OPENROUTER\_API\_KEY}
  \item эксперимент 3: \texttt{uv run python wason-2-4-6/run\_experiments.py {-}{-}models <model\_id> {-}{-}prompt-style adaptive {-}{-}reasoning-effort none {-}{-}api-key \$OPENROUTER\_API\_KEY}
\end{itemize}

\subsection{Унифицированный API-интерфейс и маршрутизация моделей}

Ключевым архитектурным решением является использование единого прокси-маршрутизатора OpenRouter (\texttt{https://openrouter.ai/api/v1}) для запросов ко всем моделям. Это устраняет артефакты, связанные с различиями в препроцессинге промптов, системных сообщениях и токенизации у разных провайдеров. Все модели запрашиваются через идентичный клиент на базе библиотеки \texttt{openai} с базовым URL OpenRouter.

Модели конфигурируются централизованно в \texttt{common\_config.py}. Список тестируемых моделей включает девять систем:
\begin{itemize}
  \item \textbf{Проприетарные:} \texttt{gpt-5.2} (OpenAI), \texttt{claude-sonnet-4.6} (Anthropic), \texttt{grok-4.1-fast} (xAI);
  \item \textbf{Открытые:} \texttt{qwen3-next-80b}, \texttt{qwen3.5-397b-a17b} (Alibaba), \texttt{deepseek-v3.2} (DeepSeek), \texttt{gemma-3-27b-it}, \texttt{gemma-3-12b-it} (Google), \texttt{glm-4.7-flash} (Zhipu AI).
\end{itemize}

Выбор моделей обеспечивает: (1) различных провайдеров, (2) различные парадигмы доступа (закрытые API и open-weight модели), (3) широкий спектр параметров (от компактных моделей 12B до масштабных систем 397B). Использование единого прокси гарантирует идентичность конфигурации запросов: формат системного сообщения, роль user/assistant и параметры токенизации остаются неизменными при переключении между моделями.

% -------------------------------------------------------
\section{Параметры запросов и механизм отказоустойчивости}
% -------------------------------------------------------

\subsection{Конфигурация параметров по экспериментам}

Параметры запросов к API дифференцируются в зависимости от методологических требований каждого эксперимента.

\textbf{Эксперимент 1 (задача Линды).} Использовались \texttt{temperature=0.0} для обеспечения максимальной детерминированности бинарного выбора, \texttt{max\_completion\_tokens=160} (достаточно для выбора опции и указания уверенности), \texttt{seed=42}. Конкурентность~--- 20 одновременных запросов, 4 попытки при ошибках, таймаут 60 секунд. Нулевая температура обусловлена требованием воспроизводимости: при повторном запросе той же виньетки модель должна вернуть идентичный ответ.

\textbf{Эксперимент 2 (задача Уэйсона).} Использовалась \texttt{temperature=1.0}, \texttt{max\_completion\_tokens=1000}, \texttt{random\_seed=42}. Единица температуры выбрана намеренно: систематические перестановки карточек (5 перестановок на правило) требуют естественной вариабельности ответов для оценки стабильности. При \texttt{temperature=0} модель могла бы выдавать идентичный ответ на все перестановки, что сделало бы невозможной оценку Consistency Rate как самостоятельной метрики. Конкурентность~--- до 60 одновременных запросов.

\textbf{Эксперимент 3 (задача 2-4-6).} Использовались \texttt{temperature=0}, \texttt{max\_turns=19}, таймаут 120 секунд на каждый многоходовой диалог, до 50 одновременных запросов. Нулевая температура обеспечивает стабильность генерации гипотез в рамках одного испытания. Максимальное число ходов (19) выбрано эмпирически: это обеспечивает достаточный горизонт для обнаружения правила, не приводя к избыточным затратам API.

\subsection{Механизм повторных попыток и repair-протокол}

Механизм повторных попыток при ошибках реализован через экспоненциальный backoff по формуле:
\begin{equation}
  \text{delay} = \min(2^{(\text{attempt}-1)}, 8) + U(0; 0{,}25) \text{ сек},
  \label{eq:backoff}
\end{equation}
где \texttt{attempt} --- номер текущей попытки (начиная с 1), а $U(0; 0{,}25)$~--- случайная добавка для предотвращения синхронизации запросов. Максимальная задержка ограничена 8 секундами.

В эксперименте 3, где требуется строго форматированный ответ (блок \texttt{```python```} с двумя строками: гипотеза и тройка), реализован дополнительный \textit{repair-протокол}. После каждого невалидного ответа формируется repair-prompt с явным указанием нарушенного требования формата:
\begin{quote}
\textit{Your previous response was invalid because: [reason]. Please respond again with EXACTLY this format: \texttt{```python H: lambda x, y, z: <hypothesis> T: [x, y, z] ```}}
\end{quote}

Выполняется до 2 дополнительных попыток ремонта. Типичные причины невалидности: выход значений за диапазон $[1, 1000]$, неполный блок (отсутствие гипотезы или тройки), использование граничных значений 0 или 1001. Все случаи нарушения формата записываются в \texttt{format\_errors.jsonl} для последующего аудита.

Наибольшее число ошибок формата зафиксировано у \texttt{grok-4.1-fast} (23 ошибки), \texttt{deepseek-v3.2} (14), \texttt{gemma-3-27b-it} (9) и \texttt{qwen3-next-80b} (7). У трёх моделей зафиксированы невалидные правила в финальном прогоне (\texttt{gemma-3-27b-it}: 2, \texttt{deepseek-v3.2}: 1, \texttt{qwen3-next-80b}: 1), исключённые из подсчёта Success Rate.

% -------------------------------------------------------
\section{Структура и генерация датасетов}
% -------------------------------------------------------

\subsection{Эксперимент 1: биографические виньетки с факториальными подстановками}

Основной датасет \texttt{linda-problem/data/conjunction\_dataset\_all.csv} содержит 130 оригинальных биографических виньеток, составленных вручную. Каждая виньетка описывает персонажа с выраженными личностными характеристиками и содержит плейсхолдеры \texttt{\{NAME\}}, \texttt{\{RACE\}} и \texttt{\{AGE\}}, заменяемые в ходе подстановки.

К каждой виньетке прилагаются пять элементов: простая опция T (профессия), коррелированное хобби F, некоррелированное хобби F\_uncorrelated и две конъюнктивные опции T\_and\_F1 и T\_and\_F2. Профессии сэмплировались из базы статистики занятости USDL/OEWS~\cite{usdl2024oews}, насчитывающей около 830 профессиональных категорий, что обеспечивает разнообразие и репрезентативность стимульного материала.

\textbf{Механизм подстановки.} Операция строковой замены (string replacement) применяется к плейсхолдерам каждой виньетки. Имена подбирались так, чтобы однозначно ассоциироваться с расовой группой согласно данным о частотности имён в США: White (Emma, James), Black (Jamal, Aisha), Hispanic (Carlos, Maria), Asian (Wei, Priya). Возрасты: 24, 38, 55 лет.

Первоначально эксперимент проводился на полной факториальной сетке: 8 имён $\times$ 3 возраста = 24 варианта подстановки на виньетку, что давало 3120 пар испытаний. После анализа результатов первых двух моделей (\texttt{glm-4.7-flash} и \texttt{qwen3.5-397b-a17b}), показавшего незначимость возрастного и proxy-gender эффектов (delta\_bias по возрастным группам в диапазоне $-0{,}005 \div +0{,}003$), сетка была сокращена до расово-значимых подстановок: 4 расовые группы $\times$ 1 имя $\times$ 1 возраст = 520 пар испытаний на модель. Данное решение, хотя и обоснованное промежуточными результатами, ограничивает возможность сравнительного анализа расового эффекта по всем девяти моделям.

\subsection{Эксперимент 2: пять контентных условий}

Пять датасетов хранятся в директории \texttt{wason-selection/input/}, каждый по 100--101 записи в формате CSV со столбцами \texttt{Rule Text} (текстовая формулировка правила) и \texttt{JSON Mapping} (отображение поверхностных значений на логические роли P, $\lnot$P, Q, $\lnot$Q).

\textbf{Протокол перестановок.} Ключевым методологическим элементом является систематическая манипуляция порядком предъявления карточек. Каждый стимул предъявлялся каждой модели в \textbf{пяти различных перестановках} четырёх позиций карточек. Дизайн следует из наблюдения Binz и Schulz~\cite{binz2023using} о том, что единственное изменение перестановки способно обратить ответ GPT-3 на каноническую задачу. Данный подход даёт \textbf{500 испытаний на модель на условие} (100 правил $\times$ 5 перестановок), позволяя оценивать распределения ответов, а не точечные оценки.

\subsection{Эксперимент 3: пятьдесят верифицированных правил}

Набор из 50 правил хранится в Google Sheets и загружается при запуске через \texttt{google\_sheets.py}. Локально сохранён манифест \texttt{wason-2-4-6/output/adaptive/none/rules\_manifest.csv} (50 строк, столбцы \texttt{code} и \texttt{category}), где \texttt{code}~--- исполняемое Python lambda-выражение вида \texttt{lambda x, y, z: x < y < z}.

Правила организованы в пять таксономических категорий: порядок и сравнение (9 правил), чётность и делимость (10), арифметические отношения (10), структурные паттерны (11), смешанные/составные (10). Все правила верифицированы на: (1) исполнимость в диапазоне $[1, 99]^3$, (2) отсутствие функциональных дубликатов (через экстенсиональную эквивалентность), (3) отсутствие тавтологий. Из 50 правил 26 удовлетворяются тройкой $(2, 4, 6)$ и 24 не удовлетворяются, что обеспечивает сбалансированную выборку для тестирования.

% -------------------------------------------------------
\section{Реализация метрик и статистического протокола}
% -------------------------------------------------------

\subsection{Метрика экстенсиональной эквивалентности (IOU)}

Ключевой методологической особенностью эксперимента 3 является оценка гипотез через экстенсиональную эквивалентность вместо текстового сравнения строк. Функция \texttt{compare\_sets(R, H, n=200000)} реализует следующий алгоритм: правило $R$ и гипотеза $H$ компилируются в Python lambda через \texttt{eval}; затем генерируется Монте-Карло-выборка из 200\,000 случайных троек $(x, y, z) \in [1, 1000]^3$ через \texttt{numpy.random.randint}; для каждой тройки вычисляются булевы векторы принадлежности, по которым считаются $|\mathbf{h}_t \cap \mathbf{r}|$, $|\mathbf{h}_t \cup \mathbf{r}|$ и итоговый IOU по формуле~\eqref{eq:iou}.

Для метрики изменения гипотез (HCR) используется отдельно кэшированная выборка из 50\,000 троек (\texttt{HYPOTHESIS\_CHANGE\_SAMPLES}), что снижает вычислительную нагрузку при сравнении пар последовательных гипотез. Использование \texttt{eval} для компиляции lambda-выражений оправдано тем, что все правила верифицированы и имеют фиксированный формат \texttt{lambda x, y, z: <expression>}, не допускающий произвольного выполнения кода.

\subsection{Метрика Severity of Bias и протокол уверенности}

В эксперименте 1 реализована непрерывная метрика силы предубеждения, основанная на самооценке уверенности модели. Функция \texttt{compute\_bias\_score(choice, confidence)} реализует следующую логику: при выборе конъюнкции (b) $\text{bias\_score} = \text{confidence} / 100$; при выборе T (a) $\text{bias\_score} = (100 - \text{confidence}) / 100$. Шкала от 0 (минимальное предубеждение) до 1 (максимальное предубеждение).

Агрегированная метрика delta\_bias вычисляется как разность средних bias\_score между correlated- и uncorrelated-условиями по формуле~\eqref{eq:delta_bias}. Эта величина отражает чистый эффект нарративной связности, очищенный от фонового уровня неуверенности модели. Парсинг уверенности реализован через regex-паттерн, извлекающий числовое значение из второй строки ответа модели.

\subsection{Статистический протокол: McNemar, Cohen's \textit{d}, FDR-коррекция}

Статистический вывод в эксперименте 1 основан на точном тесте МакНемара для согласованных пар бинарных наблюдений~\cite{mcnemar1947note}. Для каждой пары испытаний (correlated, uncorrelated) одной подставленной виньетки строилась таблица сопряжённости 2$\times$2. При $n^* = n_{12} + n_{21} > 10$ использовалась нормальная аппроксимация со статистикой $z_0 = (n_{21} - n_{12}) / \sqrt{n_{21} + n_{12}}$.

Размер эффекта оценивается через Cohen's $d = t / \sqrt{n}$, где $t$~--- значение парной t-статистики, $n$~--- число пар наблюдений. Пороги интерпретации: $0{,}2$~--- малый эффект, $0{,}5$~--- умеренный, $0{,}8$~--- большой, $>1{,}0$~--- очень большой.

При множественных сравнениях (например, между демографическими группами) применяется процедура коррекции $p$-значений Бенджамини-Хохберга (BH-FDR)~\cite{benjamini1995controlling} для контроля частоты ложных открытий на уровне $\alpha = 0{,}05$.

\subsection{Метрики дедуктивного рассуждения: EMA, MBI, CBR, Consistency Rate}

В эксперименте 2 реализованы четыре метрики. \textbf{Expected Metric Accuracy (EMA)}~--- доля испытаний с нормативно правильным набором $\{$P, $\lnot$Q$\}$. \textbf{Matching Bias Index (MBI)}~--- доля испытаний с выбором $\{$P, Q$\}$. \textbf{Confirmation Bias Rate (CBR)}~--- доля испытаний с выбором только $\{$P$\}$. \textbf{Consistency Rate}~--- доля правил, для которых модель даёт идентичный выбор карточек во всех пяти перестановках.

Все метрики вычисляются функцией \texttt{compute\_metrics.py}, которая принимает JSON-лог сырых ответов модели и отображение логических ролей карточек, возвращает агрегированные значения по каждому контентному условию и промпт-условию.

% -------------------------------------------------------
\section{Технические проблемы и их решения}
% -------------------------------------------------------

\subsection{Нарушения формата ответа в многоходовом диалоге}

Наиболее значимой технической проблемой стали систематические нарушения формата ответа в эксперименте 3. Модели должны были отвечать строго в блоке \texttt{```python```} с двумя строками: гипотеза (\texttt{H: lambda x, y, z: ...}) и тестируемая тройка (\texttt{T: [x, y, z]}). На практике модели генерировали: пояснительный текст за пределами блока, тройки с отрицательными или нулевыми значениями, тройки с числами свыше 1000, неполные блоки без гипотезы или без тройки.

Решение: двухуровневый механизм валидации. На первом уровне regex-парсер извлекает блок \texttt{```python```} и проверяет наличие обоих обязательных элементов. На втором уровне выполняется семантическая валидация: все числа тройки проверяются на принадлежность диапазону $[1, 1000]$, а гипотеза~--- на синтаксическую корректность через \texttt{compile()}. При невалидности срабатывает repair-протокол (см.~раздел~3.2.2) с указанием конкретной причины ошибки.

\subsection{Ограничение конкурентности и квоты API}

При массовом запуске эксперимента 2 (9 моделей $\times$ 5 условий $\times$ 5 перестановок $\times$ 100 правил = 22\,500 запросов) и эксперимента 3 (9 моделей $\times$ 2 промпт-условия $\times$ 50 правил $\times$ до 19 ходов = до 17\,100 запросов) возникали ограничения rate limit со стороны OpenRouter. Решение: параметрическая настройка конкурентности (от 20 до 60 одновременных запросов в зависимости от эксперимента) в сочетании с экспоненциальным backoff при получении HTTP 429.

\subsection{Недетерминированность Монте-Карло-оценки IOU}

Явная фиксация seed для генератора numpy при Монте-Карло-оценке IOU отсутствует, что приводит к незначительной недетерминированности результатов между независимыми запусками. Влияние этого эффекта оценивается как малое ввиду объёма выборки (200\,000 троек): стандартная ошибка оценки IOU при истинном значении 0{,}5 составляет $\sqrt{0{,}5 \cdot 0{,}5 / 200000} \approx 0{,}001$. Тем не менее при строгом требовании воспроизводимости рекомендуется добавить \texttt{numpy.random.seed()} перед вызовом \texttt{numpy.random.randint} в функции \texttt{compare\_sets}.

% -------------------------------------------------------
\section{Личный технический вклад автора}
% -------------------------------------------------------

Все этапы технической реализации выполнены автором самостоятельно:

\textbf{Разработано с нуля:}
\begin{itemize}
  \item Архитектура экспериментального фреймворка: асинхронный OpenRouter-клиент с механизмом retry и backoff; модульная структура с общим ядром \texttt{src/common/} и независимыми экспериментальными модулями.
  \item Система генерации и валидации датасетов: скрипт генерации биографических виньеток с факториальными подстановками; скрипт верификации правил задачи 2-4-6 на экстенсиональную эквивалентность и отсутствие тавтологий.
  \item Реализация метрик: функция Монте-Карло-оценки IOU (\texttt{compare\_sets}); вычисление Severity of Bias; агрегация McNemar, Cohen's $d$, BH-FDR; протокол перестановок карточек для Consistency Rate.
  \item Repair-протокол для эксперимента 3: двухуровневая валидация ответов с автоматической генерацией repair-промптов.
  \item Скрипты визуализации (\texttt{build\_graphs.py}) для всех трёх экспериментов с единым стилем оформления графиков.
\end{itemize}

\textbf{Использованы готовые решения:}
\begin{itemize}
  \item Библиотека \texttt{openai} --- клиент для HTTP-запросов к OpenRouter API.
  \item Библиотеки \texttt{scipy}, \texttt{scikit-learn} --- реализация статистических тестов (МакНемара, t-тест, Wilcoxon).
  \item Библиотеки \texttt{matplotlib}, \texttt{seaborn} --- рендеринг графиков.
  \item Библиотеки \texttt{gspread}, \texttt{google-auth} --- загрузка данных из Google Sheets.
  \item Менеджер пакетов \texttt{uv} --- управление зависимостями и виртуальным окружением.
\end{itemize}
